\documentclass[aspectratio=169,12pt]{beamer}
\usepackage{fontspec}
\usepackage[portuguese]{babel}
\usepackage{xcolor,listings,tikz,booktabs,amsmath,amssymb,graphicx,multicol,array,colortbl}
\usetikzlibrary{arrows.meta,positioning,shapes.geometric,calc}

\definecolor{navy}{HTML}{0D1B3E}
\definecolor{gold}{HTML}{F5A623}
\definecolor{qgreen}{HTML}{1E8B4C}
\definecolor{qblue}{HTML}{1A6EAE}
\definecolor{codebg}{HTML}{EEF3F8}

\usetheme{default}
\useinnertheme{circles}
\setbeamertemplate{navigation symbols}{}
\setbeamertemplate{footline}{%
  \hfill{\color{gray!60}\scriptsize\insertframenumber/\inserttotalframenumber}\hspace{.6em}\vspace{.4em}}

\setbeamercolor{frametitle}{bg=navy,fg=white}
\setbeamercolor{structure}{fg=navy}
\setbeamercolor{alerted text}{fg=gold}
\setbeamercolor{item}{fg=gold}
\setbeamercolor{subitem}{fg=navy}
\setbeamercolor{block title}{bg=qblue,fg=white}
\setbeamercolor{block body}{bg=qblue!8,fg=black}
\setbeamercolor{block title alerted}{bg=gold,fg=navy}
\setbeamercolor{block body alerted}{bg=gold!15,fg=navy}
\setbeamercolor{block title example}{bg=qgreen,fg=white}
\setbeamercolor{block body example}{bg=qgreen!10,fg=black}
\setbeamerfont{frametitle}{size=\large,series=\bfseries}
\setbeamerfont{block title}{series=\bfseries}

\setbeamertemplate{title page}{%
  \begin{tikzpicture}[remember picture,overlay]
    \fill[navy](current page.north west)rectangle(current page.south east);
    \fill[gold]([yshift=.7cm]current page.south west)rectangle(current page.south east);
    \node[anchor=south,yshift=.73cm,navy,font=\tiny\bfseries]at(current page.south)
      {INTENSIVÃO QUANT AI \textbullet{} IMPACT UFSCAR};
  \end{tikzpicture}
  \vspace{.65cm}
  \begin{center}
    {\color{gold}\scriptsize\bfseries INTENSIVÃO QUANT AI \textbullet{} IMPACT UFSCAR}\\[.4cm]
    {\color{white}\LARGE\bfseries\inserttitle}\\[.25cm]
    {\color{gold!85}\normalsize\insertsubtitle}\\[.5cm]
    {\color{white!70}\small\insertauthor}\\[.1cm]
    {\color{white!50}\footnotesize\insertdate}
  \end{center}}

\lstdefinestyle{python}{language=Python,
  basicstyle=\ttfamily\scriptsize\color{qblue},
  keywordstyle=\color{navy}\bfseries,
  stringstyle=\color{qgreen!80!black},
  commentstyle=\color{gray!70}\itshape,
  backgroundcolor=\color{codebg},
  frame=l,framerule=2pt,rulecolor=\color{gold},
  xleftmargin=1em,breaklines=true,showstringspaces=false,tabsize=4,
  extendedchars=false,inputencoding=ascii}
\lstset{style=python}

\author{Felipe Maldonado\\{\scriptsize VP Diretoria Quant~\textbullet{}~Impact UFSCAR}}
\date{Intensivão Quant AI 2025}
\title{Aula 08 --- Backtest Rigoroso}
\subtitle{Walk-Forward, DSR e Validação Estatística}

\begin{document}

\begin{frame}[plain]
\titlepage
\end{frame}

\begin{frame}[fragile]{Nesta Aula}
\begin{columns}[T]
  \column{0.5\textwidth}
\begin{block}{Teoria (35 min)}
\begin{itemize}
  \item Os três grandes vieses em backtests
  \item Look-ahead bias: detecção e prevenção
  \item Multiple testing bias: Harvey, Liu \& Zhu (2016)
  \item Walk-forward analysis: metodologia e implementação
  \item Deflated Sharpe Ratio (DSR): Bailey \& López de Prado (2014)
\end{itemize}
\end{block}
  \column{0.47\textwidth}
\begin{block}{Código (65 min)}
\begin{itemize}
  \item \texttt{backtest\_walkforward()} — OOS puro
  \item \texttt{calcular\_dsr()} — penalização por múltiplos testes
  \item Visualização: IS vs OOS, distribuição de retornos
  \item Salvar \texttt{retorno\_walkforward\_liquido.parquet}
\end{itemize}
\end{block}
\end{columns}
\end{frame}

\begin{frame}[fragile]{Os Três Vieses do Backtest}
\begin{columns}[T]
  \column{0.5\textwidth}
\begin{block}{1. Look-ahead Bias}
Uso de informação futura na decisão de hoje.\\
\textbf{Exemplo}: usar $r_t$ para formar o sinal de $t$.\\
\textbf{Prevenção}: sempre \texttt{shift(1)} no backtest.
\end{block}
\vspace{.3em}
\begin{block}{2. Survivorship Bias}
Só analisa ativos que sobreviveram até hoje.\\
\textbf{Exemplo}: dataset com só as empresas que não faliram.\\
\textbf{Prevenção}: incluir composição histórica do índice.
\end{block}
\vspace{.3em}
\begin{alertblock}{3. Multiple Testing Bias}
Testar $N$ variações e reportar apenas a melhor.\\
Com $N=100$ e $\alpha=5\%$, esperam-se 5 falsos positivos por acaso.\\
\textbf{Correção}: ajuste de Bonferroni ou DSR.
\end{alertblock}
  \column{0.47\textwidth}
\begin{block}{Harvey, Liu \& Zhu (2016)}

Analisaram $>$300 fatores publicados em finanças.\\
\vspace{.2em}
\textbf{Conclusão}: a maioria são \textbf{falsos positivos}.\\
O $t$-statistic mínimo aceitável evoluiu:
\begin{itemize}
  \item 1960s: $t > 2{,}0$ (Sharpe)
  \item 2000s: $t > 3{,}0$ (Harvey et al.)
  \item Hoje: DSR $> 0{,}95$ (López de Prado)
\end{itemize}
\end{block}\vspace{.3em}
\begin{exampleblock}{Walk-Forward}

Treina em janela rolling (36m), testa nos próximos 6m.
Cada retorno OOS foi gerado com parâmetros calculados \textbf{antes} do período de teste.
\end{exampleblock}
\end{columns}
\end{frame}

\begin{frame}[fragile]{Deflated Sharpe Ratio}
\begin{columns}[T]
  \column{0.5\textwidth}
\textbf{Sharpe Ratio não ajustado:}
\[ \hat{S} = \frac{\overline{r}}{\hat\sigma}\,\sqrt{12} \]

\textbf{Deflated Sharpe Ratio (Bailey \& López de Prado, 2014):}
\[ \text{DSR} = \Phi\!\left(\frac{(\hat{S} - E[\max S]) \sqrt{T-1}}{\sqrt{1 - \hat\rho\,E[\max S] + \hat{S}^2\,(\hat\kappa-1)/4}}\right) \]

Onde $E[\max S]$ cresce com o número de estratégias testadas $N$:
\[ E[\max S] \approx \sqrt{2}\,\Phi^{-1}\!\!\left(1 - \frac{1}{N}\right) \]
  \column{0.47\textwidth}
\begin{block}{Interpretação do DSR}
\begin{itemize}
  \item $\text{DSR} \in [0, 1]$: probabilidade de que o resultado é genuíno
  \item $\text{DSR} > 0{,}95$: robusto (nível de confiança 95\%)
  \item $\text{DSR} < 0{,}5$: forte suspeita de overfitting
  \item Penaliza por: número de testes, skewness e kurtosis dos retornos
\end{itemize}
\end{block}\vspace{.3em}
\begin{exampleblock}{Fórmula captura}
\begin{itemize}
  \item Correção por não-normalidade ($\hat\gamma$, $\hat\kappa$)
  \item Penalização por $N$ estratégias testadas
  \item Incerteza estatística de $\hat{S}$ com $T$ observações
\end{itemize}
\end{exampleblock}
\end{columns}
\end{frame}

\begin{frame}[fragile]{O que Vamos Construir}
\begin{columns}[T]
  \column{0.52\textwidth}
\begin{block}{Funções a implementar}
\begin{itemize}
  \item \texttt{backtest\_walkforward(sinal, ret, treino=36, teste=6)} $\to$ ret OOS
  \item \texttt{calcular\_dsr(retornos, n\_estrategias=5)} $\to$ dicionário de métricas
  \item \texttt{plot\_is\_vs\_oos(ret\_is, ret\_oos)} $\to$ comparação visual
\end{itemize}
\end{block}
  \column{0.45\textwidth}
\begin{block}{Entradas (parquets)}
\begin{itemize}
  \item \texttt{sinal\_v2.parquet}
  \item \texttt{retornos\_mensais\_limpo.parquet}
  \item \texttt{retorno\_carteira.parquet}
\end{itemize}
\end{block}
\vspace{.4em}
\begin{exampleblock}{Saídas (parquets)}
\begin{itemize}
  \item \texttt{retorno\_walkforward\_liquido.parquet}
  \item Gráfico IS vs OOS
  \item Tabela DSR e métricas OOS
\end{itemize}
\end{exampleblock}
\end{columns}
\end{frame}

\begin{frame}[fragile]{Referências}
\begin{multicols}{2}
\tiny
\begin{itemize}
  \item \textbf{Bailey, D. \& López de Prado, M.} (2014). The deflated Sharpe ratio. \textit{Journal of Portfolio Management}, 40(5).
  \item \textbf{Harvey, C., Liu, Y. \& Zhu, H.} (2016). ... and the cross-section of expected returns. \textit{Review of Financial Studies}, 29(1), 5--68.
  \item \textbf{López de Prado, M.} (2018). \textit{Advances in Financial Machine Learning}. Wiley.
  \item \textbf{White, H.} (2000). A reality check for data snooping. \textit{Econometrica}, 68(5), 1097--1126.
  \item \textbf{Romano, J. \& Wolf, M.} (2005). Stepwise multiple testing as formalized data snooping. \textit{Econometrica}, 73(4), 1237--1282.
\end{itemize}
\end{multicols}
\end{frame}


\end{document}
